% SPDX-FileCopyrightText: 2026 Bernd Zeimetz <bernd@bzed.de>
% SPDX-License-Identifier: AGPL-3.0-or-later
%
% LaTeX preamble for the PDF ("paper") rendering of IMPLEMENTATION_PLAN.md.
% Injected with pandoc --include-in-header, which in pandoc's default article
% template lands *after* fontspec/unicode-math and the verbatim machinery but
% *before* hyperref, longtable and booktabs -- hence the \AtBeginDocument
% wrappers further down.  Build with `make pdf`; see AGENTS.md section 7.

\usepackage{etoolbox}
\usepackage{xcolor}
\usepackage{framed}
\usepackage{fvextra}
\usepackage{fancyhdr}
\usepackage[explicit]{titlesec}

% ---------------------------------------------------------------- fonts ----
% The plan uses a lot of mathematical and box-drawing Unicode (l Sigma omega
% <= <=> and box drawing, block and sub/superscript characters).  No single
% well-shaped text font covers all of it, so every family gets a luaotfload
% fallback chain ending in DejaVu Sans, which covers every codepoint that
% occurs in the document, with Symbola behind it.  A glyph the chain cannot
% render produces a "Missing character" warning, which `make pdf` treats as an
% error rather than shipping a PDF with holes in it.
\directlua{
  luaotfload.add_fallback(
    "drsfallback",
    { "DejaVuSans:mode=harf;", "Symbola:mode=harf;" }
  )
}
\setmainfont{Gentium Book Plus}[RawFeature={fallback=drsfallback}]
\setsansfont{DejaVu Sans}[Scale=0.90, RawFeature={fallback=drsfallback}]
\setmonofont{DejaVu Sans Mono}[Scale=0.80, RawFeature={fallback=drsfallback}]

\linespread{1.06}

% Long identifiers set in \texttt -- metric names, config paths, API routes --
% are unbreakable words and would otherwise stick out into the margin.
% filters.lua's Code() function is where the break-after-underscore trick
% actually lives now: it hand-escapes inline code and inserts \allowbreak
% after each underscore itself, emitting a plain \texttt{...} with no macro
% override. An earlier version redefined \texttt itself (\let-ing the
% original aside and inserting a local \let\_\drsbreakableunderscore), which
% is the standard trick for this -- but \texttt is a font command the kernel
% treats as robust, or dispatches to by name in "moving" arguments (a heading
% that also has to go into the TOC, the PDF outline, and via filters.lua's
% own \markboth, the running head), and redefining it that way sent this
% kernel into unbounded recursion ("TeX capacity exceeded ... parameter stack
% size") specifically on a heading built from nothing but
% \texttt{...like_this...} -- exactly the style docs/internals/*.md's
% per-symbol headings use. Rewriting the *content* in the Lua filter, before
% it ever becomes a LaTeX command invocation, sidesteps the whole class of
% problem. Also let TeX loosen a line rather than overrun it, for the cases
% \allowbreak alone does not fully resolve; `make pdf` fails on any remaining
% overfull box wider than 2pt regardless.
\tolerance=1200
\setlength{\emergencystretch}{3em}

% ------------------------------------------------------- headings, page ----
\titleformat{\section}{\sffamily\large\bfseries}{}{0pt}{#1}
\titleformat{\subsection}{\sffamily\normalsize\bfseries}{}{0pt}{#1}
\titleformat{\subsubsection}{\sffamily\small\bfseries}{}{0pt}{#1}
% Level-4 headings ("#### 5.1.1 ...") would otherwise be run into the following
% paragraph, which hides a numbered subsection of the specification.
\titleformat{\paragraph}[block]{\sffamily\small\bfseries}{}{0pt}{#1}
\titlespacing*{\section}{0pt}{2.4ex plus 1ex minus .2ex}{1.0ex plus .2ex}
\titlespacing*{\subsection}{0pt}{1.8ex plus .8ex minus .2ex}{0.7ex plus .2ex}
\titlespacing*{\subsubsection}{0pt}{1.4ex plus .6ex minus .2ex}{0.5ex plus .2ex}
\titlespacing*{\paragraph}{0pt}{1.4ex plus .6ex minus .2ex}{0.5ex plus .2ex}

\pagestyle{fancy}
\fancyhf{}
\fancyhead[L]{\footnotesize\sffamily Proxmox Storage DRS}
\fancyhead[R]{\footnotesize\sffamily\nouppercase{\leftmark}}
\fancyfoot[C]{\footnotesize\sffamily\thepage}
\renewcommand{\headrulewidth}{0.4pt}
\renewcommand{\footrulewidth}{0pt}
% Sections are unnumbered -- the document numbers them itself ("## 5. The
% optimization problem") -- so \leftmark carries the heading text verbatim.
\renewcommand{\sectionmark}[1]{\markboth{#1}{}}

% Title page and table of contents each get pages of their own.
\pretocmd{\tableofcontents}{\clearpage}{}{}
\apptocmd{\tableofcontents}{\clearpage}{}{}

% ------------------------------------------------------------ code blocks ----
\definecolor{shadecolor}{HTML}{F4F4F1}
\fvset{fontsize=\footnotesize, breaklines=true, breakanywhere=true,
       breaksymbolleft={}, breakindent=0pt}
% Used by filters.lua for fenced blocks that carry no language, which is most
% of them: the plan's "code" is mostly formulas, PromQL fragments and ASCII
% diagrams rather than a highlightable language.  Without this they would
% render as bare, unshaded verbatim, visually unlike the highlighted few.
\newenvironment{plaincode}{\begin{snugshade}}{\end{snugshade}}

% ------------------------------------------------------------ title page ----
% The build writes \drssourcehash, \drssourcename and \drsverifycmd into a
% generated second header file (docs/.build/revision.tex for the plan, the
% equivalent for internals.pdf/pve-storage-drs-manual.pdf); the defaults here keep a bare
% `pandoc --include-in-header=header.tex` run working and describe the plan,
% since that is this project's original document. This block is shared by all
% three PDFs (documentation.md: "docs/paper/ is document-agnostic apart from
% its title block") -- \drssourcename and \drsverifycmd are what make it so:
% each build supplies its own source description and its own stamp file, and
% \drsverifycmd is always a `sha256sum --check` of that stamp rather than a
% hardcoded filename, so it reads correctly whether the source was one file
% or several concatenated.
\providecommand{\drssourcehash}{unknown}
\providecommand{\drssourcename}{IMPLEMENTATION\_PLAN.md}
\providecommand{\drsverifycmd}{sha256sum --check docs/IMPLEMENTATION\_PLAN.pdf.sha256}

\makeatletter
\renewcommand{\maketitle}{%
  \begin{titlepage}%
    \centering
    \vspace*{0.22\textheight}
    {\sffamily\bfseries\Huge\@title\par}%
    \vspace{1.6em}
    \rule{0.55\textwidth}{0.4pt}\par
    \vspace{1.6em}
    {\large\@author\par}%
    \vspace{0.8em}
    {\@date\par}%
    \vfill
    {\footnotesize
     Rendered from \texttt{\drssourcename},
     SHA-256 \texttt{\drssourcehash}\\
     (\texttt{\drsverifycmd} in the repository must
     succeed).\\[0.9em]
     Copyright \textcopyright{} 2026 Bernd Zeimetz.
     Licensed under the GNU Affero General Public License,
     version 3 or later.\par}
    \vspace{1.2em}
  \end{titlepage}%
  % article's titlepage resets the counter to 1, which would make the printed
  % page numbers run one behind the PDF's.  Count the title page instead.
  \setcounter{page}{2}%
}
\makeatother

% ------------------------------------------- deferred: loaded after this ----
\AtBeginDocument{%
  % Section 11 and 15 hold wide tables; shrink them and let cells wrap rather
  % than overflow into the margin.
  \AtBeginEnvironment{longtable}{\footnotesize}%
  \AtBeginEnvironment{tabular}{\footnotesize}%
  \setlength{\LTpre}{0.8\baselineskip}%
  \setlength{\LTpost}{0.8\baselineskip}%
  \hypersetup{%
    bookmarksnumbered=false,
    bookmarksopen=true,
    bookmarksopenlevel=1,
    pdfdisplaydoctitle=true,
  }%
  \urlstyle{same}%
}
